% GENERATED, edit the French unit or the translation table
\mxnewpage\section*{Solutions EV1-T2}\label{EV1-T2-CORR}
\textbf{MCQ:} 1 D, 2 C, 3 B, 4 D, 5 A, 6 A, 7 C, 8 A, 9 B, 10 C, 11 D, 12 C. Each correct answer: \num{0,5} point; total 6.
\paragraph{Checks of the fluency questions.} Order of operations; sum of fractions; subtraction of exponents; multiplying coefficient \num{1,2}; division by \num{0,8}; product \num{1,25} times \num{0,8}; solving an affine equation; vertical increment divided by horizontal; mean; angle divided by 360; complement to 1; substitution in brackets.
\subsection*{Exercise 1 -- Two growth processes}
\begin{enumerate}
\item A arithmetic, B geometric; $a_2=50,b_2=120$. \textit{Marking: 1 point. Result and basic justification expected; half a point possible for a correct but unfinished approach.}
\item A: $\dfrac{4(30+60)}{2}=180$; B: $30(1+2+4+8)=450$. \textit{Marking: 2 points. Method or justification: 1; calculation and conclusion: 1.}
\item $b_1=60>a_1=40$. For $n\geq1$, the increment of B is $30\times2^n\geq60$, greater than 10; each step increases the gap. \textit{Marking: 2 points. Method or justification: 1; calculation and conclusion: 1.}
\end{enumerate}
\subsection*{Exercise 2 -- Imposed roots}
\begin{enumerate}
\item $f=a(x+1)(x-5)$; $-5a=10$, so $a=-2$. \textit{Marking: 1 point. Result and basic justification expected; half a point possible for a correct but unfinished approach.}
\item Negative coefficient; f positive between the roots, negative elsewhere. Solutions $\intof{-\infty}{-1}\cup\intfo{5}{+\infty}$. \textit{Marking: 2 points. Method or justification: 1; calculation and conclusion: 1.}
\item $f=-2x^2+8x+10=-2(x-2)^2+18$. Maximum 18 at 2. \textit{Marking: 2 points. Method or justification: 1; calculation and conclusion: 1.}
\end{enumerate}
\subsection*{Exercise 3 -- A controlled approximation}
\begin{enumerate}
\item Rate $6+h$. \textit{Marking: 1 point. Result and basic justification expected; half a point possible for a correct but unfinished approach.}
\item $f\prime(3)=6$, $f(3)=10$; $y=6x-8$. \textit{Marking: 1 point. Result and basic justification expected; half a point possible for a correct but unfinished approach.}
\item Approximation \num{10,12}; exact $\num{10,1204}$. Exact minus approximate $\num{0,0004}=h^2$. \textit{Marking: 2 points. Method or justification: 1; calculation and conclusion: 1.}
\end{enumerate}
\paragraph{Calibration.} Target duration: 20--25 min for the fluency questions, 80--85 min for the exercises, 10--15 min for proofreading. These durations are editorial estimates to be measured in class. Accept any correct method; do not penalise twice a numerical error that is propagated consistently. The proposed marking scheme is not a decision of the team to harmonise marking.
